% Intellectual Property Assignment Agreement. Content from IPA.docx.
\documentclass{paridyum}
\begin{document}

\paridyumletterhead
\paridyumdateref{\dkDocDate}{\dkDocRef}
\paridyumtitle{Intellectual Property Assignment Agreement}

This Intellectual Property Assignment Agreement (``Agreement'') is entered into
between \textbf{\dkOrgName} (``Paridyum'') and \textbf{\dkPIdentityLegalName} (``Intern'').

\section*{1. Purpose}
During the internship, the Intern may create or contribute to software, designs,
documentation, research, ideas, or other materials for Paridyum. This Agreement
defines ownership of that work.

\section*{2. Work Product}
``Work Product'' means work created by the Intern specifically as part of their
assigned internship responsibilities or Paridyum projects. This may include:
\begin{itemize}[leftmargin=*, noitemsep]
    \item source code and software;
    \item scripts and APIs;
    \item algorithms;
    \item databases and schemas;
    \item technical documentation;
    \item product designs and prototypes;
    \item UI/UX designs;
    \item graphics and other creative materials;
    \item research and technical specifications;
    \item test cases and automation;
    \item presentations and written materials;
    \item inventions, improvements, and discoveries; and
    \item other materials created specifically for Paridyum.
\end{itemize}

\section*{3. Ownership and Assignment}
To the extent permitted by applicable law, the Intern assigns to Paridyum all
transferable rights, title, and interest in Work Product created specifically for
Paridyum during the internship. This includes the right to use, reproduce, modify,
distribute, publish, license, commercialize, and create derivative works from the
Work Product. The assignment is intended to apply worldwide and for the full duration
of the applicable intellectual-property rights, to the extent permitted by law.

\section*{4. Intern's Pre-Existing Work}
The Intern continues to own intellectual property created independently before the
internship or outside the scope of Paridyum work. If the Intern wants to use their
existing code, libraries, designs, tools, or other intellectual property in a Paridyum
project, they must disclose it to Paridyum before incorporating it.

\vspace{0.4em}
\textbf{Pre-existing intellectual property to be used:} \dkPPreExistingIp

\section*{5. Third-Party Materials}
The Intern must not knowingly include third-party proprietary or confidential material
in Paridyum projects without authorization. This includes material belonging to:
\begin{itemize}[leftmargin=*, noitemsep]
    \item previous employers;
    \item other companies;
    \item clients;
    \item universities;
    \item other internships; or
    \item other individuals.
\end{itemize}

\section*{6. Open-Source Software}
The Intern may use open-source software where appropriate, provided that:
\begin{itemize}[leftmargin=*, noitemsep]
    \item its license permits the intended use;
    \item applicable license requirements are followed; and
    \item any significant licensing restrictions are disclosed to Paridyum.
\end{itemize}
The Intern should obtain approval before introducing software with licensing terms
that may require Paridyum's proprietary source code to be disclosed or distributed.

\section*{7. Further Documentation}
The Intern agrees to reasonably assist Paridyum if additional signatures or documents
are required to confirm or record ownership of Work Product created during the
internship. This obligation may continue after the internship ends where reasonably
necessary to confirm rights arising from work completed during the internship.

\section*{8. Future Paridyum Entity}
If Paridyum later operates through a registered company or other legal entity, rights
obtained under this Agreement may be assigned or transferred to that entity where
permitted by applicable law. The Intern agrees to reasonably cooperate with any
documentation required to confirm such transfer.

\section*{9. Confidentiality}
Work Product that has not been publicly released by Paridyum must be treated as
confidential. The Intern's confidentiality obligations are governed by the separate
\textbf{Non-Disclosure Agreement}.

\section*{10. Portfolio Use}
The Intern may state on their resume, CV, or professional profile that they worked on
projects at Paridyum. However, the Intern may not publish or display proprietary source
code, confidential documentation, internal screenshots, unreleased designs, or other
confidential Work Product without prior written approval from Paridyum.

\section*{11. Intern's Confirmation}
The Intern confirms that, to the best of their knowledge:
\begin{itemize}[leftmargin=*, noitemsep]
    \item their original contributions do not knowingly violate another person's
          intellectual-property rights;
    \item they will not knowingly copy restricted proprietary material into Paridyum
          projects;
    \item they will disclose relevant third-party materials or licenses; and
    \item they will disclose relevant pre-existing intellectual property before
          incorporating it into Paridyum Work Product.
\end{itemize}

\section*{12. Governing Law}
This Agreement will be governed by the laws of \dkOrgJurisdictionCountry. Subject to
applicable law, disputes relating to this Agreement will be subject to the appropriate
courts in \textbf{\dkOrgJurisdictionForum}.

\section*{13. Acceptance}
By signing below, the parties confirm that they have read, understood, and agreed to
the terms of this Intellectual Property Assignment Agreement.

\vspace{1.5em}
\paridyumofficers{\dkSigCeoName}{\dkSigCeoDesignation}{\dkSigCtoName}{\dkSigCtoDesignation}

\vspace{1.5em}
\paridyumcountersign{Intern}{\dkPIdentityLegalName}

\end{document}
